\section{Formulation of the Model Problem}
\label{sec:model_problem}

{\color{red} write intro to section}

Let $\Omega \subset \mathbb{R}^2$ denote the domain representing the
terrain, and let $a,b \in \Omega$ denote the prescribed start and end
points, respectively. An admissible path is represented by a
sufficiently smooth curve
\[
\gamma : [0,1] \rightarrow \mathbb{R}^2,
\]
satisfying the boundary conditions
\[
\gamma(0)=a,
\qquad
\gamma(1)=b.
\]
The collection of admissible curves satisfying these boundary
conditions, together with any additional constraints imposed by the
terrain, is denoted by $\Gamma$.

The terrain is described by a cost field
\[
c:\Omega\rightarrow\mathbb{R},
\]
where $c(x,y)$ represents the cost associated with traversing the
terrain at the point $(x,y)$. In the simplest formulation, the cost of
a path is given by
\begin{equation}
J_{\mathrm{terrain}}(\gamma)
=
\int_\gamma c(x,y)\,ds.
\label{eq:terrain_functional}
\end{equation}
This functional accounts for both the cost of the terrain encountered
by the path and the distance travelled through each region.

In order to incorporate additional considerations relevant to the
path-planning problem, penalties for curvature and total path length
are introduced. The resulting functional is written in the compact
form
\begin{equation}
\boxed{
J(\gamma)
=
\int_\gamma
\left[
w_c c(x,y)
+
\lambda\kappa^2
+
\mu
\right]\,ds
}
\label{eq:ModelProblem}
\end{equation}
where $w_c$, $\lambda$, and $\mu$ are non-negative parameters controlling
the relative contribution of the terrain cost, curvature, and length,
respectively.

The first term in the integrand,
\[
w_c c(x,y),
\]
represents the cost associated with traversing the terrain. The
parameter $w_c$ allows the contribution of the terrain cost to be
scaled relative to the additional penalties.

The second term,
\[
\lambda\kappa^2,
\]
penalises curvature. In particular, the square of the curvature assigns
increasingly greater cost to regions of the path with large curvature,
thereby discouraging sharp changes in direction and favouring smoother
paths.

Finally, the constant term $\mu$ represents a penalty on the total
length of the path. Since
\[
L(\gamma)=\int_\gamma 1\,ds,
\]
the contribution of this term to the functional is
\[
\mu\int_\gamma 1\,ds
=
\mu L(\gamma).
\]
Consequently, the functional in \eqref{eq:ModelProblem} can equivalently
be written as
\begin{equation}
J(\gamma)
=
w_c\int_\gamma c(x,y)\,ds
+
\lambda\int_\gamma\kappa^2\,ds
+
\mu L(\gamma).
\label{eq:expanded_model}
\end{equation}

The corresponding variational problem is therefore
\begin{equation}
\inf_{\gamma\in\Gamma} J(\gamma).
\label{eq:extended_problem}
\end{equation}

The three terms in \eqref{eq:ModelProblem} introduce a trade-off between
terrain cost, path smoothness, and path length. Increasing $w_c$
favours paths which avoid regions of high terrain cost, while increasing
$\lambda$ places greater emphasis on smoothness. Similarly, increasing
$\mu$ makes longer paths increasingly expensive and therefore favours
shorter routes. The numerical investigation that follows seeks to
understand how these competing terms influence the resulting
approximate minimisers.

\subsection{Arc Length and Path Length}
\label{subsec:arc_length}

Let the curve $\gamma$ be parametrised by
\[
\gamma(t)=(x(t),y(t)),
\qquad t\in[0,1].
\]
The length of the curve is given by
\begin{equation}
L(\gamma)
=
\int_0^1
\|\gamma'(t)\|\,dt.
\label{eq:path_length}
\end{equation}

More generally, the arc length measured from $\gamma(0)$ to
$\gamma(t)$ is
\begin{equation}
s(t)
=
\int_0^t
\|\gamma'(\tau)\|\,d\tau.
\label{eq:arc_length}
\end{equation}

Thus, an integral with respect to arc length can be expressed in terms
of the parameter $t$ according to
\begin{equation}
\int_\gamma f(x,y)\,ds
=
\int_0^1
f(\gamma(t))
\|\gamma'(t)\|\,dt.
\label{eq:line_integral}
\end{equation}

In particular, the terrain component of the functional can be written
as
\begin{equation}
\int_\gamma c(x,y)\,ds
=
\int_0^1
c(\gamma(t))
\|\gamma'(t)\|\,dt.
\end{equation}

\subsection{Curvature Penalty}
\label{subsec:curvature_penalty}

The curvature of a curve measures the rate at which its direction
changes with respect to arc length. When $\gamma$ is parametrised by
arc length $s$, the tangent vector satisfies
\[
\left\|\frac{d\gamma}{ds}\right\|=1,
\]
and the curvature is given by
\begin{equation}
\kappa(s)
=
\left\|
\frac{d^2\gamma}{ds^2}
\right\|.
\label{eq:curvature_arclength}
\end{equation}

The curvature contribution to the functional is consequently
\begin{equation}
J_{\mathrm{curvature}}(\gamma)
=
\int_\gamma \kappa^2\,ds.
\label{eq:curvature_functional}
\end{equation}

The use of $\kappa^2$ ensures that curvature is penalised independently
of its direction, while placing a greater penalty on regions where the
path bends sharply. This provides a mechanism for controlling the
smoothness of the resulting path.

The formulation in \eqref{eq:ModelProblem} therefore provides a
mathematical representation of the competing requirements of the
path-planning problem. In the following sections, this continuous
variational problem is discretised and implemented computationally,
allowing approximate minimisers to be obtained using simulated
annealing.

\section{Euler-Lagrange Equations}

{\color{red} This is just to dot down some maths before discussiuoins are proper write up}

First investigate the problem in terms of $s$ instead of $x,y$ we can rewrite the \eqref{eq:ModelProblem} as 
\begin{equation}
    \label{eq:ModelPrblem_s}
    J[\gamma] = \int_{s_0}^{s_1} c(\gamma(s)) + \lambda|\gamma''(s)|^2 + \mu \, ds.
\end{equation}

Using the Second order Euler Ladgrange Equation {\color{red} referance this}
\[
\frac{\partial F}{\partial\gamma} - \frac{d}{ds}\left[\frac{\partial F}{\partial\gamma'}\right] + \frac{d^2}{d^2s}\left[\frac{\partial F}{\partial\gamma''}\right] = 0
\]

we are left with 
\begin{equation}
    \label{eq:ModelProblem_EL_Solution}
    \gamma^{(4)}(s) = -\frac{1}{2\lambda} \nabla_{\gamma} c(\gamma(s))
\end{equation}

{\color{red} I'm a bit lost on how to solve this or what we can get form this}

{\color{red} I did have a go with a hill function for c}

to mimic the python code let $c(s)$ be {\color{red} I will re order so the python code mimics this}
\begin{equation}
    \label{eq:hill_C}
    c(\gamma(s)) = He^{-\frac{|\gamma(s) - p|^2}{2\sigma^2}}
\end{equation} 
Where $H$ is hil magnitude (as in the differance between max cost and backgound cost, $p$ is centrepoint of the hill, and $\sigma$ is the spead of the hill,
\[
H = c_{max} - c_0 
\]
\[
p=(s_c) = (x_c,y_c)
\]
subing this $c$ into \eqref{eq:ModelProblem_EL_Solution} gives
\[
    \gamma^{(4)}(s) = \frac{H}{2\lambda\sigma^2} e^{-\frac{|\gamma(s) - p|^2}{2\sigma^2}}(\gamma(s) - p)
\]

Since this is a 4th order non linear ODE we can see that typciall analyticval methods cant help us and a different method is required.